\section{The Legacy of the Byronic Hero}

The Romantic Era was one that idealized an unspecified past as more natural and where people's lives were less fragmented, slower paced, and fuller. Some assert that the movement was something of a revolt against the emphasis on the rationalism of the Enlightenment, but more troubling was its attack on traditional morality. Indeed, it seems odd that a movement that idealized the past due to its simplicity and contentment would attack the standard that made such things possible.

A longing for this fuller life is a recurring characteristic of Lord Byron's protagonists. Characters with these traits have evolved into an archetype of their known: Byronic heroes. For the young and dissatisfied, these anti-heroes have much appeal. Highly contradictory beings, these figures have imposed themselves to self-exile, as their restlessness and shame condemns them to always look to a new horizon where they believe that a better life will await them.

The Byronic Hero is amoral. While genuinely searching for a pure life, the heroes instability causes him to take action that violates the most basic codes of conduct. As a strong, charismatic, and attractive man, he is able to accomplish great deeds, gain fame, and seduce the innocent, rural women he idealizes, but his achievements ultimately prove fleeting.

These figures challenge the establishment and have a contempt for hierarchy. When their destruction is taken to its limit, they break the social bonds that tie society together, but some realize their error before this happens.

The traits of these figures were embodied in 60s `Beat' or `Hipster' counterculture. Jack Kerouac, Allen Ginsberg, and other writers of the time were real people who brought the sense of dissatisfaction and revolt against the `normal' life to its logical extent. Like Romantic contempt for the artificiality of Enlightenment thought and society, the Beats saw in American bourgeois life nothing except paralysis. To them, the way of life offered represented nothing more than a death before dying, and they put themselves to the task of rediscovering `life.'

The Beats took the restlessness represented by the Byronic Hero and brought it to its most extreme extent. Destroying every bond and limit became a total, uncompromising mission. Freedom came to be understood as a temporary, overload of the senses; it was to them something that could only be grasped for a short time and had to be continuously rediscovered through ruthless destruction of the self and others. Normalcy came to interpreted as oppression. Purity, here, becomes inverted in its totality. One does not find happiness through innocence, chastity, and contentment, instead, restlessness is cured by an attitude of complete apathy towards the future and the development of an ability to impulsively erupt at any given moment.

The utter success of this social revolution lay in the lack of comprehension that those who claim to be its opponents have regards to its success. A new, contra-civilizational standard was set. Restlessness and apathy have become the new norm, and the social institutions have been reformed not to push people back toward social ways of life, but rather to normalize the abnormal and to create dissatisfaction and confusion where there had not been before. Rene Guenon explains this process perfectly in the Reign of Quantity and the Sign of the Times, describing the transformation of the institutions as the ``Counter-initiation.''

For the individual to overcome these processes, delusion about their nature cannot be permitted. Buddha derived a formula to remind oneself of when such conditions arise: ``No, I do not want this. This does not belong to me. This, I am not. This is not myself.'' This renunciation of the world is meant to be no arbitrary thing. Rather, it is a refusal to mix oneself with the ignoble and corrupt oneself in the process. Of course, the Christian teaching is similar: To be in the world but not of the world. Also, ``Don't copy the behavior and customs of this world, but let God transform you into a new person by changing the way you think. Then you will learn to know God's will for you, which is good and pleasing and perfect.''

These renunciations are not meant to mean withdraw from the world or denial of the will, although in plenty of mistaken interpretations of these religions that is exactly what is practiced. Rather, one should will, but not in a self-deluded matter. Seeking to ``obtain positions of power'' in order to ``reform the system from within'' is not a working proposition without first a more general spiritual realization. It was our own restlessness and misplaced desire that brought us to the crisis we are in now. Until the next horizon is understood to be the reconquering of the self, and not more chasing delusions, we will get nowhere, and our fate is already sealed.


\flrightit{Posted on 2011-02-02 by VisionsOfGlory14}

\begin{center}* * *\end{center}

\begin{footnotesize}
\begin{sffamily}
\texttt{Count Cagliostro on 2011-02-05 at 07:34 said:}

Excellent text!

PS Do you happen to know the name of the creator of the painting?

\hfill

\texttt{VisionsOfGlory14 on 2011-02-05 at 09:57 said:}

It's a painting of Faust and Mephistopheles. Can't tell you who the painter is.

\hfill

\texttt{Helmholtz on 2011-02-12 at 04:36 said:}

Apparently that's an Arnold Böcklin work, also a self-portrait:
\textit{http://en.wikipedia.org/wiki/Arnold\_Böcklin}

\hfill

\end{sffamily}
\end{footnotesize}

